\documentclass[10pt,a4paper]{amsart}
\usepackage[margin=19mm]{geometry}
\usepackage{amsmath,amssymb,mathtools}
\usepackage[scr=rsfs,cal=euler]{mathalfa}
\usepackage{xfrac}
\usepackage[unicode,hidelinks,bookmarks=false]{hyperref}
\newcommand{\ie}{i.e.\ }
\newcommand{\cf}{cf.\ }
\newcommand{\resp}{respectively\ }
\newcommand{\Subsection}{\subsection}
\providecommand{\nobreakdash}{\nobreak\hbox{-}}
\setlength{\emergencystretch}{2em}
\multlinegap0pt
\usepackage{amssymb}
\usepackage{enumerate}
\usepackage{cancel}
\usepackage[shortlabels]{enumitem}
\newtheorem{theorem}{Theorem}
\newtheorem{lemma}{Lemma}
\newcommand{\Ho}{\mathcal H}
\newcommand{\Tor}{\mathrm{Tor}}
\newcommand{\V}{\mathcal V}
\title{Sub-Laplacian generalized curvature dimension inequalities on Riemannian foliations}
\author{Fabrice Baudoin, Guang Yang}
\date{Source: arXiv:2601.02035v1 (CC BY 4.0)}
\begin{document}
\maketitle

\noindent\emph{Source and licence.} Sub-Laplacian generalized curvature dimension inequalities on Riemannian foliations. Version arXiv:2601.02035v1 (CC BY 4.0), distributed by arXiv under the Creative Commons Attribution 4.0 International licence, http://creativecommons.org/licenses/by/4.0/, as recorded in the arXiv licence metadata for this exact version and shown on the abstract page https://arxiv.org/abs/2601.02035v1. Full licence notice: \texttt{licenses/arxiv-2601.02035-CC-BY-4.0.txt}.\par\smallskip

\noindent\textit{Editorial scope.} Counted unit: the article's theorem Bochner (source0.tex lines 607--620). The article's complete original proof of it is delivered with it (source0.tex lines 783--855, one \texttt{proof} environment of 73 lines, which the article places away from the statement). No mathematics is written, repaired or re-derived: the statement and the proof are the article's own lines, in the article's own notation, and no retained line is substituted or re-flowed.  Retained uncounted as the source-local context the proof needs: lines 625--629; lines 650--659; lines 703--709; lines 737--745.  Nothing was omitted from any retained span, so the package carries no omission record.  No retained line contains a citation, so the wrapper carries no bibliography and no bibliography entry.  No retained line contains a figure environment, an image-inclusion command or drawing code, so no image is delivered and the package declares no asset dependency.  Editorial formatting only, in addition: an AMS \texttt{amsart} preamble that restates the article's own declarations and macros used by the retained lines, namely the environments \texttt{theorem}, \texttt{lemma} on separate counters, together with the standard packages those lines need.  The retained environments print in this document as Theorem 1, Lemma 1 in order of appearance, whereas the article prints the same items with the numbering of its own full document.  The complete native source of the article as distributed in the arXiv e-print for this exact version is bundled unchanged as \texttt{sources/192131/source0.tex}.
	\begin{lemma}\label{symmetrized-hessian}
		\begin{align*}
			|\nabla_\Ho \nabla_\Ho f|^2 =| \mathrm{Hess}_\Ho^{\nabla, \mathrm{sym}} f|^2 +\frac{1}{4}( J_{\nabla_\V f} , J_{\nabla_\V f} )_\Ho
		\end{align*}
	\end{lemma}

	\begin{lemma}\label{partial hessian} 
		If $X_1,\cdots, X_n$ is a local horizontal orthonormal frame then
		\[
		\sum_{i=1}^n \langle[\nabla_{\nabla_\mathcal{H}f} \nabla_{X_i}-\nabla_{\nabla_{\nabla_\mathcal{H}f} X_i}] \nabla f, X_i \rangle=\langle \nabla_{\mathcal{H}} f , \nabla_{\mathcal{H}}\overset{\circ}\Delta_{\mathcal{H}} f \rangle
		\]
		and
		\[
		\sum_{i=1}^n \langle[\nabla_{\nabla_\mathcal{V}f} \nabla_{X_i}-\nabla_{\nabla_{\nabla_\mathcal{V}f} X_i}] \nabla f, X_i \rangle=\langle \nabla_{\mathcal{V}} f , \nabla_{\mathcal{V}}\overset{\circ}\Delta_{\mathcal{H}} f \rangle.
		\]
	\end{lemma}

	\begin{lemma}\label{horizontal-part-grad}
		We have
		\begin{align*}
			\nabla_{\mathcal{H}}u_1&=\nabla_{\nabla_\mathcal{H}f} (\nabla_\mathcal{H}f)+(J_{\nabla f}\nabla_\mathcal{H}f)_{\mathcal{H}},\\
			\nabla_{\mathcal{H}}u_2&=\nabla_{\nabla_\mathcal{V}f} (\nabla_\mathcal{H}f)+(J_{\nabla  f}\nabla_\mathcal{V}f)_{\mathcal{H}}.
		\end{align*}
	\end{lemma}

    \begin{lemma}\label{mean curvature part}
\begin{align*}
			\frac{1}{2}\mathrm H | \nabla_\Ho f|^2-\langle \nabla_{\mathcal{H}} f , \nabla_{\mathcal{H}} \mathrm H f \rangle=- \nabla^{\mathrm{sym}}\mathrm H(\nabla_\Ho f,\nabla_\Ho f)-\left\langle \mathrm{Tor}^\nabla (\mathrm H, \nabla_\Ho f), \nabla f \right\rangle 
		\end{align*}
		and
		\begin{align*}
			\frac{1}{2}\mathrm H | \nabla_\V f|^2-\langle \nabla_{\mathcal{V}} f , \nabla_{\mathcal{V}}\mathrm H f \rangle=-2\nabla^{\mathrm{sym}}\mathrm H(\nabla_\V f,\nabla_\Ho f)-\left\langle \mathrm{Tor}^\nabla (\mathrm H, \nabla_\V f), \nabla f \right\rangle.
		\end{align*}
    \end{lemma}

	\begin{theorem}\label{Bochner}
		Let $f \in C^\infty (M)$ and  $X_1,\cdots,X_n$ be a local orthonormal frame of horizontal vector fields. We have
		\begin{align*}
			\frac{1}{2}\Delta_{\mathcal{H}} | \nabla_\Ho f|^2=&\langle \nabla_{\mathcal{H}} f , \nabla_{\mathcal{H}}\Delta_{\mathcal{H}} f \rangle+2\sum_i\langle \nabla_{X_i} \nabla_\V f, \mathrm{Tor}^\nabla(\nabla_{\mathcal{H}}f, X_i) \rangle+| \mathrm{Hess}_\Ho^{\nabla, \mathrm{sym}} f|^2\\
			&- \langle \nabla_\V f, \delta_\Ho \mathrm{Tor}^\nabla ( \nabla_\Ho f ) \rangle+\mathrm{Ric}_\Ho^\nabla(\nabla_{\mathcal{H}} f, \nabla_{\mathcal{H}} f )-(\mathrm{Tor}^\nabla (\nabla_\Ho f),\mathrm{Tor}^\nabla (\nabla_\V f) )_\Ho\\
			&+\frac{1}{4}( J_{\nabla_\V f} , J_{\nabla_\V f} )_\Ho+ \nabla^{\mathrm{sym}} \mathrm H(\nabla_\Ho f,\nabla_\Ho f)+\left\langle \mathrm{Tor}^\nabla (\mathrm H, \nabla_\Ho f), \nabla f \right\rangle -\tau(\nabla_\Ho f,\nabla_\Ho f)
		\end{align*}
		and
		\begin{align*}
			\frac{1}{2}\Delta_{\mathcal{H}} | \nabla_\V f|^2=&\langle \nabla_{\mathcal{V}} f , \nabla_{\mathcal{V}}\Delta_{\mathcal{H}} f \rangle+2\sum_i\langle \nabla_{X_i} \nabla f, \mathrm{Tor}^\nabla(\nabla_{\mathcal{V}}f, X_i) \rangle+\mid \nabla_{\mathcal{H}} \nabla_{\mathcal{V}} f \mid^2\\ 
			&-\langle \nabla f, \delta_\Ho \mathrm{Tor}^\nabla ( \nabla_\V f ) \rangle-(\mathrm{Tor}^\nabla (\nabla_\V f),\mathrm{Tor}^\nabla (\nabla_\V f) )_\Ho +\mathrm{Ric}^\nabla_\Ho(\nabla_\V f,\nabla_\Ho f)\\
            &+2\nabla^{\mathrm{sym}}\mathrm H(\nabla_\Ho f,\nabla_\V f)+\left\langle \mathrm{Tor}^\nabla (\mathrm H, \nabla_\V f), \nabla f \right\rangle-\tau(\nabla_\V f,\nabla_\Ho f).
		\end{align*}
	\end{theorem}

	\begin{proof}[Proof of Theorem \ref{Bochner}]
		
		Plug in our result from Lemma \ref{horizontal-part-grad}, we get that for every $X \in \mathfrak{X}(\Ho)$,
		\begin{align*}
			&\langle \nabla_X \nabla_{\mathcal{H}}u_1 ,X \rangle\\
			=&\langle \nabla_X \nabla_{\nabla_\mathcal{H}f} \nabla_{\mathcal{H}}f, X \rangle+\langle \nabla_X (J_{\nabla f}\nabla_\mathcal{H}f)_{\mathcal{H}}, X \rangle\\
			=&\langle \nabla_X \nabla_{\nabla_\mathcal{H}f} \nabla f, X \rangle+\langle \nabla_X (J_{\nabla f}\nabla_\mathcal{H}f), X \rangle\\
			=&\langle \nabla_{\nabla_\mathcal{H}f} \nabla_X  \nabla f, X \rangle+\langle \mathrm{Riem}^{\nabla}(X, \nabla_{\mathcal{H}} f) \nabla f, X\rangle+{\langle \nabla_{[X, \nabla_\mathcal{H}f]}  \nabla f, X \rangle}+\langle \nabla_X (J_{\nabla f}\nabla_\mathcal{H}f), X \rangle\\
			=&\langle \nabla_{\nabla_\mathcal{H}f} \nabla_X  \nabla f, X \rangle+\langle \mathrm{Riem}^{\nabla}(X, \nabla_{\mathcal{H}} f) \nabla_\Ho f, X\rangle+{\langle \nabla_{\nabla_X \nabla_\mathcal{H}f-\nabla_{\nabla_\mathcal{H}f} X-\mathrm{Tor}^{\nabla}(X, \nabla_\mathcal{H}f) } \nabla f, X \rangle}\\
			+&\langle \nabla_X (J_{\nabla f}\nabla_\mathcal{H}f), X \rangle\\
			=&\langle[\nabla_{\nabla_\mathcal{H}f} \nabla_X-\nabla_{\nabla_{\nabla_\mathcal{H}f} X}] \nabla f, X \rangle+ \langle\mathrm{Riem}^{\nabla}(X, \nabla_{\mathcal{H}} f) \nabla_\Ho f, X\rangle+\langle \nabla_{\nabla_X \nabla_\mathcal{H}f-\Tor^{\nabla}(X, \nabla_\mathcal{H}f) } \nabla f, X \rangle\\
			+&\langle \nabla_X (J_{\nabla f}\nabla_\mathcal{H}f), X \rangle.
		\end{align*}
		Recall that we have for all $A,B\in \mathfrak{X}(TM)$,
		\[ \langle \nabla_A \nabla f, B \rangle=\langle \nabla_B \nabla f, A \rangle+\langle \Tor^\nabla(B,A), \nabla f \rangle.   \]
		Thus
		\begin{align*}
			&\langle \nabla_{\nabla_X \nabla_\mathcal{H}f} \nabla f, X \rangle-\langle \nabla_{\Tor^{\nabla}(X, \nabla_\mathcal{H}f) } \nabla f, X \rangle\\
			=&\langle \nabla_{X} \nabla f, \nabla_X \nabla_\mathcal{H}f  \rangle+\langle \Tor^{\nabla}(X, \nabla_X \nabla_\mathcal{H}f ), \nabla f \rangle\\
			-&\langle \nabla_X\nabla f , \Tor^\nabla(X, \nabla_{\mathcal{H}}f) \rangle-\langle \Tor^\nabla(X, \Tor^{\nabla}(X, \nabla_\mathcal{H}f) ), \nabla f \rangle.
		\end{align*}
		Moreover, since $\nabla$ is metric-compatible, we have
		\begin{align*}
			&\langle \nabla_X (J_{\nabla f}\nabla_\mathcal{H}f), X \rangle\\
			=&X\langle J_{\nabla f}\nabla_\mathcal{H}f, X \rangle-\langle J_{\nabla f}\nabla_\mathcal{H}f, \nabla_X X \rangle  \\
			=&X\langle \nabla f, \Tor^\nabla(\nabla_\mathcal{H}f, X) \rangle-\langle   \nabla f, \Tor^\nabla(\nabla_{\mathcal{H}}f, \nabla_X X)\rangle\\
			=&\langle \nabla_X \nabla f, \Tor^\nabla(\nabla_{\mathcal{H}}f, X) \rangle+\langle \nabla f, \nabla_X (\Tor^\nabla(\nabla_{\mathcal{H}}f, X)) \rangle-\langle \nabla f,  \Tor^\nabla(\nabla_{\mathcal{H}}f, \nabla_X X) \rangle \\
			=&\langle \nabla_X \nabla f, \Tor^\nabla(\nabla_{\mathcal{H}}f, X) \rangle+\langle \nabla f, (\nabla_X \Tor^\nabla)(\nabla_{\mathcal{H}}f, X) \rangle+\langle \nabla f,  \Tor^\nabla(\nabla_X \nabla_{\mathcal{H}}f,  X) \rangle.
		\end{align*}
		
		Therefore
		\begin{align*}
			& \langle \nabla_{\nabla_X \nabla_\mathcal{H}f-\Tor^{\nabla}(X, \nabla_\mathcal{H}f) } \nabla f, X \rangle\
			+\langle \nabla_X (J_{\nabla f}\nabla_\mathcal{H}f), X \rangle \\
			=&\langle \nabla_{X} \nabla f, \nabla_X \nabla_\mathcal{H}f  \rangle+2\langle \nabla_X \nabla f, \Tor^\nabla(\nabla_{\mathcal{H}}f, X) \rangle+\langle \nabla f, (\nabla_X \Tor^\nabla)(\nabla_{\mathcal{H}}f, X) \rangle \\
			&-\langle \Tor^\nabla(X, \Tor^{\nabla}(X, \nabla_\mathcal{H}f) ), \nabla f \rangle \\
			=&|\nabla_X \nabla_\mathcal{H}f  |^2+2\langle \nabla_X \nabla_\V f, \Tor^\nabla(\nabla_{\mathcal{H}}f, X) \rangle+\langle \nabla_\V f, (\nabla_X \Tor^\nabla)(\nabla_{\mathcal{H}}f, X) \rangle \\
			&-\langle \Tor^{\nabla}(X, \nabla_\mathcal{V}f) ,\Tor^{\nabla}(X, \nabla_\mathcal{H}f)\rangle-\langle \Tor^\nabla(X, \Tor^{\nabla}(X, \nabla_\mathcal{H}f) ), \nabla_\Ho f \rangle.
		\end{align*}
		In the last line, we used the fact that the torsion of the connection $\nabla$ satisfies for $U,V \in\mathfrak{X}(\V)$ and $X,Y \in \mathfrak{X}(\Ho)$
		\[
		\langle \Tor^\nabla (U,X),V \rangle=\langle \Tor^\nabla (V,X),U \rangle, \qquad \Tor^\nabla (X,Y) \in \mathcal{X}(\V).
		\]
		We put everything together and get		
		\begin{align*}
			\langle \nabla_X \nabla_{\mathcal{H}}u_1 ,X \rangle=&\langle[\nabla_{\nabla_\mathcal{H}f} \nabla_X-\nabla_{\nabla_{\nabla_\mathcal{H}f} X}] \nabla f, X \rangle +|\nabla_X \nabla_\mathcal{H}f  |^2+\langle\mathrm{Riem}^{\nabla}(X, \nabla_{\mathcal{H}} f) \nabla_\Ho f, X\rangle \\
			&+2\langle \nabla_X \nabla_\V f, \Tor^\nabla(\nabla_{\mathcal{H}}f, X) \rangle+\langle \nabla_\V f, (\nabla_X \Tor^\nabla)(\nabla_{\mathcal{H}}f, X) \rangle \\
			& -\langle \Tor^{\nabla}(X, \nabla_\mathcal{V}f) ,\Tor^{\nabla}(X, \nabla_\mathcal{H}f)\rangle-\langle \Tor^\nabla(X, \Tor^{\nabla}(X, \nabla_\mathcal{H}f) ), \nabla_\Ho f \rangle.
		\end{align*}
		Finally, let $X$ range over a horizontal frame, and we get the result  from Lemmas \ref{mean curvature part}, \ref{symmetrized-hessian} and \ref{partial hessian} after summing up. 
        
        
        We now turn to the second Bochner's formula in the vertical directions. The computation follows the same lines, we therefore only show the main steps.  Similar to the previous case, we have
		\begin{align*}
			&\langle \nabla_X \nabla_{\mathcal{H}}u_2 ,X \rangle\\
			=&\langle[\nabla_{\nabla_\mathcal{V}f} \nabla_X-\nabla_{\nabla_{\nabla_\mathcal{V}f} X}] (\nabla f), X \rangle+\langle\mathrm{Riem}^{\nabla}(X, \nabla_{\mathcal{V}} f) \nabla_\Ho f, X\rangle+\langle \nabla_{\nabla_X \nabla_\mathcal{V}f-\Tor^{\nabla}(X, \nabla_\mathcal{V}f) } (\nabla f), X \rangle\\
			&+\langle \nabla_X (J_{\nabla f}\nabla_\mathcal{V}f)_{\mathcal{H}}, X \rangle.
		\end{align*} 
		Then, as before, we have
		\begin{align*}
			& \langle \nabla_{\nabla_X \nabla_\mathcal{V}f-\Tor^{\nabla}(X, \nabla_\mathcal{V}f) } \nabla f, X \rangle\
			+\langle \nabla_X (J_{\nabla_\V f}\nabla_\mathcal{V}f), X \rangle \\
			=&|\nabla_X \nabla_\mathcal{V}f  |^2+2\langle \nabla_X \nabla f, \Tor^\nabla(\nabla_{\mathcal{V}}f, X) \rangle+\langle \nabla f, (\nabla_X \Tor^\nabla)(\nabla_{\mathcal{V}}f, X) \rangle \\
			&-\langle \Tor^{\nabla}(X, \nabla_\mathcal{V}f) ,\Tor^{\nabla}(X, \nabla_\mathcal{V}f)\rangle-\langle \Tor^\nabla(X, \Tor^{\nabla}(X, \nabla_\mathcal{V}f) ), \nabla_\Ho f \rangle.
		\end{align*}
		Therefore, we obtain
		\begin{align*}
			\langle \nabla_X \nabla_{\mathcal{H}}u_2 ,X \rangle=&\langle[\nabla_{\nabla_\mathcal{V}f} \nabla_X-\nabla_{\nabla_{\nabla_\mathcal{V}f} X}] \nabla f, X \rangle +|\nabla_X \nabla_\mathcal{V}f  |^2+2\langle \nabla_X \nabla f, \Tor^\nabla(\nabla_{\mathcal{V}}f, X) \rangle \\
			&+\langle \nabla f, (\nabla_X \Tor^\nabla)(\nabla_{\mathcal{V}}f, X) \rangle  -\langle \Tor^{\nabla}(X, \nabla_\mathcal{V}f) ,\Tor^{\nabla}(X, \nabla_\mathcal{V}f)\rangle \\
             &-\langle \Tor^\nabla(X, \Tor^{\nabla}(X, \nabla_\mathcal{V}f) ), \nabla_\Ho f \rangle+\langle\mathrm{Riem}^{\nabla}(X, \nabla_{\mathcal{V}} f) \nabla_\Ho f, X\rangle
		\end{align*}
		and the conclusion follows from Lemmas \ref{mean curvature part} and then  \ref{partial hessian} after summing up over a local horizontal orthonormal frame.
	\end{proof}

\end{document}
